
\documentclass[12pt]{amsart}

\usepackage[margin=1in]{geometry}
\usepackage{amsmath,amssymb,amsthm}

\newtheorem{theorem}{Theorem}[section]

\newcommand{\Z}{\mathbb{Z}}
\newcommand{\R}{\mathbb{R}}
\newcommand{\e}{\mathrm{e}}
\newcommand{\C}{\mathrm{C}}
\newcommand{\SL}{\mathrm{SL}}


\begin{document}

\section{Definition and Notation}
We fix some notation for studying an obtuse regime away from the boundary.

\subsection*{Truncated obtuse region}
Fix $\eta\in(0,1/6)$. For $n\ge 1$, define the set of integer pairs
\begin{equation}\label{eq:Hn}
\mathcal H_n(\eta):=\{(p,q)\in\mathbb Z^2:\ \eta n\le p,q,\ \ p+q<\tfrac{n}{2},\ \gcd(p,q,n)=1\}.
\end{equation}

\subsection*{Largest prime divisor}
Fix notation as follows. Let $P:=P^+(n)$ be the largest prime divisor of $n$, and write
\[
n=P^{\alpha}m,\qquad \alpha:=v_P(n)\ge 1,\qquad \gcd(P,m)=1.
\]
We also let
\begin{equation}
U_n:=\{a\in\{1,2,\dots,n\}:\ \gcd(a,n)=1\}
\end{equation}
be the reduced residue system modulo $n$.

\subsection*{Least nonnegative residue}
Write $[x]_n$ for the least nonnegative residue of $x$ modulo $n$.

\subsection*{A counting function}
We define the indicator of the interval $I_m:=\{1,2,\dots,m\}\subset\Z/n\Z$ by
\[
1_{I_m}(x):=
\begin{cases}
1,& 1\le [x]_n\le m,\\
0,& \text{otherwise.}
\end{cases}
\]
For $(p,q)\in\mathcal{H}_n(\eta)$, we define
\[
m_p:=2p-1\qquad and \qquad m_q:=2q-1,
\]
and the counting function (using indicator functions)
\begin{equation}
\label{eq:Sdef}
S(p,q):=\sum_{a\in U_n} 1_{I_{m_p}}(ap)\,1_{I_{m_q}}(aq).
\end{equation}

\section{Main Theorem}


\begin{theorem}[Analytic engine: a lower bound for $S(p,q)$]\label{thm:engine}
If $\eta\in(0,1/6)$ and $\theta\in(0,1)$, then we have
$$
\lim_{\substack{n\rightarrow +\infty\\
P^+(n)\geq n^{\theta}}} \frac{|\left \{ (p,q) \in \mathcal{H}_n(\eta) \ : \ \gcd\!\big(q,P^+(n)\big)=1  \ \ and \ \ S(p,q)<5\right\}|}{|\mathcal{H}_n(\eta)|}=0.
$$
\end{theorem}


\end{document}